\chapter{Conclusions}
\label{chap:conclusions}

This thesis set out to narrow the marketing-quality gap between large brands and the independent retailers that compete with them. The starting observation was that recent multimodal models can already produce near-photographic visuals from a short prompt, yet the step from a model's raw output to a flyer that a small shop is willing to print or post remains stubborn. The argument running through the chapters has been that this step is closed by the protocol that sits around the model, not by a bigger model.

Every objective set out in Section~\ref{sec:objectives} has been met. The platform that grew out of this work is a single application a retailer can use from their phone or from their desk: they sign in once, assemble a product library either by photographing items themselves or by searching a curated catalogue from a sample shot (Figure~\ref{fig:result-references}), and from then on receive a fresh marketing idea every morning, tuned to the weather, to upcoming holidays and to events near the shop. One tap turns an idea into a catalogue, an announcement or a wallpaper. Another tap exports it as a paper-sized PDF ready for the printer (Figure~\ref{fig:result-print}). The retailer never sees a prompt box or a model selector at any point in this loop. Real outputs from each of the paths are collected in Section~\ref{sec:sample-assets}, and the whole thing runs as a regular production service rather than a controlled experiment.

The hardest single problem was product fidelity. Generative models do not reliably reproduce a specific product from a reference photograph, and no retailer is going to publish a flyer in which their own product is subtly wrong. Chapter~\ref{chap:contributions} describes the hybrid pipeline that addresses this: the model designs the surrounding scene with reserved placeholder regions, the compositor pastes the retailer's actual photographs into those regions, and a judge then checks the candidate against a short rubric (no residual outlines, no visible seams, no obvious distortion). The pipeline retries up to three times when the judge rejects, which lifts the effective acceptance rate from roughly $70\%$ on a single shot to roughly $97\%$ at the cap (Table~\ref{tab:hyb-acceptance}). The cap is deliberately tunable: a fourth or fifth trial would push the rate to $\approx 99.2\%$ or $\approx 99.8\%$, but every retailer would pay the extra latency and the extra model call, even though only the worst few per cent of runs would benefit. Three trials proved a sensible balance for the small-retailer workflow, and the table makes it straightforward to revisit that choice if the operational data later justifies it.

What the platform leaves the retailer with is the practical thing the thesis set out to deliver: a tool that turns the shop's own photographs, brand and daily context into marketing assets the retailer is willing to put their name on, without ever asking the retailer to think about prompts, models or pipelines. Section~\ref{sec:sample-assets} shows real outputs across every surface the platform offers, from the daily idea card the retailer reads in the morning to the print-ready PDF that ends up on a flyer in the window. The reading is the one the thesis argued for from the first chapter on: the step from a general-purpose model to an output an independent retailer will actually publish is closed by what sits around the model.

\section{Future improvements}
\label{sec:future-work}

Several extensions would meaningfully strengthen the platform.

\begin{itemize}
    \item Self-serve billing via Stripe: replace the admin-grant flow with Stripe Checkout for one-off credit packs and Stripe Billing for subscriptions, reconciling charges against the existing ledger through webhooks.
    \item Engagement feedback loop: pull per-post reach, engagement and click-through from the connected social channels and feed them back into the recommendation engine as a per-shop signal, so future ideas favour patterns that have worked for that retailer.
    \item Per-shop writer adaptation: use periodic LoRA refreshes on each shop's accepted-versus-rejected ideas to let the writer converge to the shop's voice without retraining the base model; the hot-swappable adapter infrastructure is already in place.
    \item Short-form video generation: extend the pipeline beyond stills to $5$ to $15$\,s vertical clips for Reels, TikTok and Shorts; the hybrid contract translates naturally, with the compositor placing real product photos onto each generated frame.
    \item Controlled benchmark for the recommendation engine: run a rigorous study with a power-sized frozen evaluation set, blind human raters from the target audience, and a per-failure-mode breakdown, in place of the current LLM-as-judge and small pairwise pass.
\end{itemize}
